\chapter{正方体与长方体}\label{ch:g5:solids}

告别平面：\omterm{def:g5:solids:def}{立体图形}是现实世界中的形状——
骰子、麦片盒、罐头、球。本章学习给它们命名，
数一数它们的面和棱，把它们展开成\omterm{def:g5:solids:net}{展开图}，并
通过数小\omterm{def:g5:solids:def}{正方体}来测量它们的\omterm{def:g5:solids:volume}{体积}。（\omterm{def:g5:solids:volume}{体积}公式见
\cref{ch:g6:measure}）

\section{立体图形家族}

\begin{definition}[立体图形及其相关词汇]\label{def:g5:solids:def}
占据空间的物体称为\emph{立体图形}\index{立体图形}。平的侧面称为
\emph{面}\index{面}，面与面相交于\emph{棱}\index{棱}；
棱与棱相交于\emph{顶点}\index{顶点}。
这一家族的成员如下：
\begin{itemize}
  \item \emph{正方体}\index{正方体}：$6$ 个正方形的面，$12$ 条棱，
        $8$ 个顶点；
  \item \emph{长方体}\index{长方体}：$6$ 个长方形的面，
        $12$ 条棱，$8$ 个顶点；
  \item \emph{圆柱}\index{圆柱}：两个圆面和一个
        侧面（如罐头）；
  \item \emph{棱锥}\index{棱锥}：一个\omterm{def:g4:geometry:polygon}{多边形}底面和若干三角形，它们交于
        同一个顶点；
  \item \emph{球}：一个完全圆滑的曲面，
        没有面，没有棱。
\end{itemize}
\end{definition}

\begin{omfigure}
\begin{tikzpicture}[scale=0.62]
  % cube
  \draw[thick] (0,0) rectangle (1.8,1.8);
  \draw[thick] (0,1.8) -- (0.7,2.5) -- (2.5,2.5) -- (1.8,1.8);
  \draw[thick] (2.5,2.5) -- (2.5,0.7) -- (1.8,0);
  \node[below, font=\small] at (1.2,-0.3) {正方体};
  % box
  \begin{scope}[xshift=4.4cm]
  \draw[thick] (0,0) rectangle (2.8,1.4);
  \draw[thick] (0,1.4) -- (0.6,2.0) -- (3.4,2.0) -- (2.8,1.4);
  \draw[thick] (3.4,2.0) -- (3.4,0.6) -- (2.8,0);
  \node[below, font=\small] at (1.7,-0.3) {长方体};
  \end{scope}
  % cylinder
  \begin{scope}[xshift=9.6cm]
  \draw[thick] (0,0.35) ellipse (0.85 and 0.3);
  \draw[thick] (-0.85,0.35) -- (-0.85,2.2);
  \draw[thick] (0.85,0.35) -- (0.85,2.2);
  \draw[thick] (0,2.2) ellipse (0.85 and 0.3);
  \node[below, font=\small] at (0,-0.3) {圆柱};
  \end{scope}
  % pyramid
  \begin{scope}[xshift=12.6cm]
  \draw[thick] (-0.9,0.2) -- (0.9,0.2) -- (1.35,0.75) -- cycle;
  \draw[thick] (-0.9,0.2) -- (0.15,2.4) -- (0.9,0.2);
  \draw[thick] (1.35,0.75) -- (0.15,2.4);
  \draw[dashed] (-0.9,0.2) -- (-0.45,0.75) -- (1.35,0.75);
  \draw[dashed] (-0.45,0.75) -- (0.15,2.4);
  \node[below, font=\small] at (0.2,-0.3) {棱锥};
  \end{scope}
  % ball
  \begin{scope}[xshift=16.4cm]
  \draw[thick] (0,1.2) circle (1.1);
  \draw[dashed] (0,1.2) ellipse (1.1 and 0.35);
  \node[below, font=\small] at (0,-0.3) {球};
  \end{scope}
\end{tikzpicture}

{\small \omterm{def:g5:solids:def}{立体图形}展览。虚线表示看不见的棱——
即\omterm{def:g5:solids:def}{立体图形}背面的棱，仿佛透过\omterm{def:g5:solids:def}{立体图形}看到。}
\end{omfigure}

\begin{example}[数一数正方体]\label{ex:g5:solids:count}
数一数骰子：$6$ 个面（标有数字 $1$ 到 $6$），$8$ 个
\omterm{def:g5:solids:def}{顶点}（即角），$12$ 条棱（用手指沿着棱划一划：上面 $4$ 条，
下面 $4$ 条，竖着的 $4$ 条）。还有
$6 + 8 = 12 + 2$——这一奇妙规律对\omterm{def:g5:solids:def}{长方体}和
\omterm{def:g5:solids:def}{棱锥}也成立，见\cref{ex:g8:solids:count}。
\end{example}

\section{展开图}

\begin{definition}[展开图]\label{def:g5:solids:net}
\omterm{def:g5:solids:def}{立体图形}的\emph{展开图}\index{展开图}是其所有面的平面画法，
各面沿棱相连，可折叠成该\omterm{def:g5:solids:def}{立体图形}——
即把\omterm{def:g5:solids:def}{立体图形}像回收前压平的纸箱一样拆开。
\end{definition}

\begin{omfigure}
\begin{tikzpicture}[scale=0.5]
  \draw[gray!40, thin] (0,0) grid (4,3);
  \draw[thick, omDef] (0,1) rectangle (1,2);
  \draw[thick, omDef] (1,1) rectangle (2,2);
  \draw[thick, omDef] (2,1) rectangle (3,2);
  \draw[thick, omDef] (3,1) rectangle (4,2);
  \draw[thick, omDef] (1,2) rectangle (2,3);
  \draw[thick, omDef] (1,0) rectangle (2,1);
  \node[below, font=\small] at (2,-0.5) {正方体的展开图};
  \begin{scope}[xshift=7.5cm]
  \draw[gray!40, thin] (0,0) grid (4,3);
  \draw[thick, omThm] (0,1) rectangle (1,2);
  \draw[thick, omThm] (1,1) rectangle (2,2);
  \draw[thick, omThm] (2,1) rectangle (3,2);
  \draw[thick, omThm] (3,1) rectangle (4,2);
  \draw[thick, omThm] (0,2) rectangle (1,3);
  \draw[thick, omThm] (0,0) rectangle (1,1);
  \node[below, font=\small] at (2,-0.5) {另一种正确的展开图};
  \end{scope}
  \begin{scope}[xshift=15cm]
  \draw[gray!40, thin] (0,0) grid (4,3);
  \draw[thick, omExo] (0,1) rectangle (1,2);
  \draw[thick, omExo] (1,1) rectangle (2,2);
  \draw[thick, omExo] (2,1) rectangle (3,2);
  \draw[thick, omExo] (3,1) rectangle (4,2);
  \draw[thick, omExo] (0,2) rectangle (1,3);
  \draw[thick, omExo] (1,2) rectangle (2,3);
  \node[below, font=\small] at (2,-0.5) {不是展开图！（为什么？）};
  \end{scope}
\end{tikzpicture}

{\small 六个正方形的三种排列。前两种可以折成
\omterm{def:g5:solids:def}{正方体}；第三种不能——折叠时，有两个正方形落在同一个面上，
\omterm{def:g5:solids:def}{正方体}便无法封闭。剪下来试一试！}
\end{omfigure}

\begin{example}[读展开图]\label{ex:g5:solids:read}
骰子上，相对两面的点数和为 $7$。在\omterm{def:g5:solids:net}{展开图}上，相对的面是指
一行中隔开恰好一个正方形的两个面（或绕着一个角隔开的两个面），
而不是相邻的面！在折叠之前，在\omterm{def:g5:solids:net}{展开图}上标出这些相对的面，
是锻炼空间想象的好办法。
\end{example}

\section{体积：数小正方体}

\begin{definition}[用计数法求体积]\label{def:g5:solids:volume}
由单位\omterm{def:g5:solids:def}{正方体}搭成的\omterm{def:g5:solids:def}{立体图形}的\emph{体积}\index{体积}，
就是它所含\omterm{def:g5:solids:def}{正方体}的个数。单位可以是厘米\omterm{def:g5:solids:def}{正方体}
（cm$^3$）或米\omterm{def:g5:solids:def}{正方体}（m$^3$）。
\end{definition}

\begin{example}[逐层数长方体]\label{ex:g5:solids:layers}
一个\omterm{def:g5:solids:def}{长方体}长 $4$ 个、宽 $3$ 个、高 $2$ 个小\omterm{def:g5:solids:def}{正方体}，底层
有 $4 \times 3 = 12$ 个小\omterm{def:g5:solids:def}{正方体}，共 $2$ 层：
\[
12 \times 2 = 24 \text{ 个小正方体} .
\]
逐层计数总是行得通的，它还悄悄证明了
\cref{ch:g6:measure}中的公式 $V = L \times w \times h$。
\end{example}

\begin{omfigure}
\begin{tikzpicture}[scale=0.52]
  \foreach \y in {0,1} {
    \foreach \x in {0,...,3} {
      \foreach \z in {0,1,2} {
        \begin{scope}[shift={({\x+0.4*\z},{\y+0.28*\z})}]
          \draw[fill=omDef!12] (0,0) rectangle (1,1);
          \draw[fill=omDef!22] (0,1) -- (0.4,1.28) -- (1.4,1.28)
            -- (1,1) -- cycle;
          \draw[fill=omDef!30] (1,0) -- (1.4,0.28) -- (1.4,1.28)
            -- (1,1) -- cycle;
        \end{scope}
      }
    }
  }
  \node[below, font=\small] at (2.6,-0.4)
    {每层 $4 \times 3$ 个，$2$ 层：$24$ 个小正方体};
\end{tikzpicture}

{\small 由单位小\omterm{def:g5:solids:def}{正方体}搭成的 $4 \times 3 \times 2$ \omterm{def:g5:solids:def}{长方体}，
逐层计数。}
\end{omfigure}

\begin{example}[同样的小正方体，不同的形状]\label{ex:g5:solids:rearrange}
八个单位小\omterm{def:g5:solids:def}{正方体}，可以搭成一个 $2 \times 2 \times 2$ 的\omterm{def:g5:solids:def}{正方体}、
一根 $8 \times 1 \times 1$ 的长条，或一块 $4 \times 2 \times 1$ 的薄板，
共三种不同的\omterm{def:g5:solids:def}{立体图形}，\omterm{def:g5:solids:volume}{体积}都是 $8$ 个小\omterm{def:g5:solids:def}{正方体}。
\omterm{def:g5:solids:volume}{体积}衡量的是材料的多少，而不是形状。
\end{example}

\section{练习}

\begin{exercise}[$\star$]\label{exo:g5:solids:1}
请说出生活中形状像下列\omterm{def:g5:solids:def}{立体图形}的实物：\omterm{def:g5:solids:def}{正方体}；\omterm{def:g5:solids:def}{长方体}；\omterm{def:g5:solids:def}{圆柱}；
球；\omterm{def:g5:solids:def}{棱锥}。
\end{exercise}

\begin{exercise}[$\star$]\label{exo:g5:solids:2}
\omterm{def:g5:solids:def}{长方体}有多少个面、多少条棱、多少个\omterm{def:g5:solids:def}{顶点}？请验证这个规律：
面数 $+$ \omterm{def:g5:solids:def}{顶点}数 $=$ 棱数 $+ 2$。
\end{exercise}

\begin{exercise}[$\star$]\label{exo:g5:solids:3}
\omterm{def:g5:solids:def}{长方体}的哪些面是完全相同的？（把六个面两两分组。）
\omterm{def:g5:solids:def}{正方体}的面有什么特殊之处？
\end{exercise}

\begin{exercise}[$\star$]\label{exo:g5:solids:4}
画出棱长为 $2$ 厘米的\omterm{def:g5:solids:def}{正方体}的\omterm{def:g5:solids:net}{展开图}（可用本章的十字形模型），
剪下来并折叠成\omterm{def:g5:solids:def}{正方体}。
\end{exercise}

\begin{exercise}[$\star$]\label{exo:g5:solids:5}
在骰子的十字形\omterm{def:g5:solids:net}{展开图}上，填入数字 $1$ 到 $6$，使相对两面
的点数之和为 $7$。（先用\cref{ex:g5:solids:read}找出
相对的面对。）
\end{exercise}

\begin{exercise}[$\star$]\label{exo:g5:solids:6}
数一数单位小\omterm{def:g5:solids:def}{正方体}的个数：长 $5$、宽 $2$、高 $3$ 的\omterm{def:g5:solids:def}{长方体}；棱长为 $3$ 的
\omterm{def:g5:solids:def}{正方体}；一个 L 形立体，由一块 $3 \times 2 \times 1$ 的薄板
再在上面叠放一块 $1 \times 2 \times 1$ 的薄板组成。
\end{exercise}

\begin{exercise}[$\star$]\label{exo:g5:solids:7}
一个\omterm{def:g5:solids:def}{长方体}的某一层恰好含 $12$ 个单位小\omterm{def:g5:solids:def}{正方体}。请给出两种可能的
每层形状（$L \times w$），并对每种形状求出该\omterm{def:g5:solids:def}{长方体}有
$3$ 层时的\omterm{def:g5:solids:volume}{体积}。
\end{exercise}

\begin{exercise}[$\star$]\label{exo:g5:solids:8}
用 $27$ 个单位小\omterm{def:g5:solids:def}{正方体}，能搭出棱长是多少的\omterm{def:g5:solids:def}{正方体}？用 $64$ 个呢？
（想一想：把逐层计数倒过来想。）
\end{exercise}

\begin{exercise}[$\star\star$]\label{exo:g5:solids:9}
棱长为 $1$ 厘米的方糖，装在一个长 $10$ 厘米、宽 $5$ 厘米、
高 $4$ 厘米（内部尺寸）的盒子里。这个盒子能装下多少块方糖？
如果一家人每天用 $6$ 块方糖，
那么一盒能用多少天（请合理地\omterm{def:g4:numbers:round}{四舍五入}）？
\end{exercise}

\begin{exercise}[$\star\star$]\label{exo:g5:solids:10}
一个 $3 \times 3 \times 3$ 的\omterm{def:g5:solids:def}{正方体}，外表面涂成红色，然后
切成 $27$ 个单位小\omterm{def:g5:solids:def}{正方体}。恰有
$3$ 个面涂有颜色的小\omterm{def:g5:solids:def}{正方体}有几个？恰有 $2$ 个面的呢？恰有 $1$ 个面的呢？一个面都没有的呢？（检验：这四个数之和应为
$27$。）
\end{exercise}
